\begin{tikzpicture}[
  x=1mm,y=1mm,
  font=\figurefont\scriptsize,
  panel/.style={draw=BookRule,fill=NNPanel,line width=.85pt,rounded corners=1.5mm},
  axis/.style={knowledge arrow,draw=BookMuted,line width=.8pt},
  viewone/.style args={#1}{circle,fill=#1,draw=#1,minimum size=3.2mm,inner sep=0pt},
  viewtwo/.style args={#1}{circle,fill=white,draw=#1,line width=1pt,minimum size=3.8mm,inner sep=0pt},
  samplelabel/.style={font=\figurefont\scriptsize,text=BookInk},
  note/.style={font=\figurefont\scriptsize,text=BookMuted,align=center}
]
  % 所有面板用同一映射：(z_1,z_2) -> (中心横坐标+8z_1,37+8z_2) mm。
  % 坐标范围均为[-2,2]^2；点位直接使用矩阵中的教学数值。
  \draw[panel] (.5,.5) rectangle (51.5,66.5);
  \draw[panel] (58.5,.5) rectangle (110.5,66.5);
  \draw[panel] (117.5,.5) rectangle (168.5,66.5);
  \foreach \cx in {26,84.5,143}{
    \begin{scope}[shift={(\cx,37)},x=8mm,y=8mm]
      \draw[axis] (-2,0) -- (2,0) node[right] {$z_1$};
      \draw[axis] (0,-2) -- (0,2) node[above] {$z_2$};
      \draw[BookMuted,line width=.6pt] (1,-.04) -- (1,.04)
        node[below=.4mm] {$1$};
      \draw[BookMuted,line width=.6pt] (-.04,1) -- (.04,1)
        node[left=.4mm] {$1$};
    \end{scope}
  }
  \node[font=\figurefont\bfseries\footnotesize,text=BookInk] at (26,62) {(a) 常数表示};
  \node[viewtwo=BookAmber,minimum size=7mm] at (26,37) {};
  \node[viewtwo=NNHidden,minimum size=5.8mm] at (26,37) {};
  \node[viewtwo=BookTeal,minimum size=4.6mm] at (26,37) {};
  \node[viewone=BookBlue] at (26,37) {};
  \node[samplelabel,anchor=west] at (30,40) {$A=B=C=D$};
  \node[note] at (26,12) {$\bm{Z}=\bm{Z}'=\begin{bmatrix}0&0\\0&0\\0&0\\0&0\end{bmatrix}$\\配对误差为$0$，各维方差为$0$};

  \node[font=\figurefont\bfseries\footnotesize,text=BookInk] at (84.5,62) {(b) 单方向变化};
  \begin{scope}[shift={(84.5,37)},x=8mm,y=8mm]
    \draw[BookRule,line width=.6pt,dashed] (-1.7,-1.7) -- (1.7,1.7);
    \foreach \x/\y/\sampleid/\samplecolor in {-1.5/-1.5/A/BookBlue,-.5/-.5/B/BookTeal,.5/.5/C/NNHidden,1.5/1.5/D/BookAmber}{
      \node[viewtwo=\samplecolor] at ({\x+.08},{\y-.05}) {};
      \node[viewone=\samplecolor] at (\x,\y) {};
      \node[samplelabel,anchor=south east,xshift=-1mm,yshift=1mm] at (\x,\y) {\sampleid};
    }
  \end{scope}
  \node[note] at (84.5,10.5) {$\bm{Z}\approx\bm{Z}'=\begin{bmatrix}-1.5&-1.5\\-.5&-.5\\.5&.5\\1.5&1.5\end{bmatrix}$\\样本有方差，但两个坐标重复同一变化};

  \node[font=\figurefont\bfseries\footnotesize,text=BookInk] at (143,62) {(c) 示意维度中分散};
  \begin{scope}[shift={(143,37)},x=8mm,y=8mm]
    \foreach \x/\y/\sampleid/\samplecolor in {-1/-1/A/BookBlue,-1/1/B/BookTeal,1/-1/C/NNHidden,1/1/D/BookAmber}{
      \node[viewtwo=\samplecolor] at ({\x+.08},{\y-.05}) {};
      \node[viewone=\samplecolor] at (\x,\y) {};
      \node[samplelabel,anchor=south east,xshift=-1mm,yshift=1mm] at (\x,\y) {\sampleid};
    }
  \end{scope}
  \node[note] at (143,10.5) {$\bm{Z}\approx\bm{Z}'=\begin{bmatrix}-1&-1\\-1&1\\1&-1\\1&1\end{bmatrix}$\\两维有方差且示例协方差为$0$};

  \node[viewone=BookInk] at (5,-5) {};
  \node[note,anchor=west] at (8,-5) {视图A};
  \node[viewtwo=BookInk] at (37,-5) {};
  \node[note,anchor=west] at (40,-5) {对应视图B};
  \node[note,anchor=west,text=FigureLoss] at (76,-5) {教学构造；分散不等于形成有用类别};
\end{tikzpicture}
